\chapter{Data model}
\label{chap:data}

Postgres holds users, GitHub installations, repositories, reviews, findings, and the code index. Figure~\ref{fig:er} is the main graph. An arrow points from a row to the row it references. \code{reviews} also references \code{repositories} and \code{users}; those two links are omitted from the drawing so the pull-request chain stays readable. \code{webhook\_deliveries} has no foreign key. Eval tables (\code{eval\_runs}, \code{eval\_case\_results}) hang off \code{users} and are omitted from the figure for the same reason.

\begin{figure}[H]
  \centering
  \resizebox{\linewidth}{!}{% Main tables. Arrows point from a row to the row it references.
\begin{tikzpicture}[x=1cm,y=1cm]
  \node[entity] (users) at (7.6,7.15) {users};

  \node[entity] (keys) at (1.55,5.35) {api keys};
  \node[entity] (llm) at (4.55,5.35) {llm settings};
  \node[entity] (notif) at (7.55,5.35) {notifications};
  \node[entity] (inst) at (10.55,5.35) {installations};
  \node[entity] (app) at (13.55,5.35) {github app};

  \node[entity] (chunks) at (7.55,3.55) {code chunks};
  \node[entity] (repos) at (10.55,3.55) {repositories};
  \node[entity] (settings) at (13.55,3.55) {repo settings};

  \node[entity] (hooks) at (1.55,1.7) {webhook\\deliveries};
  \node[entity] (prs) at (10.55,1.7) {pull requests};

  \node[entity] (events) at (7.55,0.15) {review events};
  \node[entity] (reviews) at (10.55,0.15) {reviews};
  \node[entity] (findings) at (13.55,0.15) {findings};

  \draw[rel] (keys) -- (users);
  \draw[rel] (llm) -- (users);
  \draw[rel] (notif) -- (users);
  \draw[rel] (inst) -- (users);
  \draw[rel] (app) -- (users);
  \draw[rel] (repos) -- (inst);
  \draw[rel] (settings) -- (repos);
  \draw[rel] (chunks) -- (repos);
  \draw[rel] (prs) -- (repos);
  \draw[rel] (reviews) -- (prs);
  \draw[rel] (events) -- (reviews);
  \draw[rel] (findings) -- (reviews);

  \node[entityplain, anchor=west] (note) at (3.55,1.7) {no foreign key};
\end{tikzpicture}
}
  \caption{Primary relationships. \code{webhook deliveries} records a GitHub delivery id so a retry is not processed twice.}
  \label{fig:er}
\end{figure}

\section{Tables}

\begin{table}[H]
  \centering
  \caption{What each table is for.}
  \label{tab:tables}
  \small
  \begin{tabular}{@{}lp{8.8cm}@{}}
    \toprule
    Table & Role \\
    \midrule
    \code{users} & Clerk user id and email. \\
    \code{installations} & A GitHub App installation. \code{user\_id} is null until Connect GitHub finishes, because the webhook can arrive first. \\
    \code{repositories} & One connected repo: enabled flag, index status, last indexed commit, and the embedding model the index was built with. \\
    \code{repo\_settings} & Strictness, ignored paths, custom rules, and the maximum number of comments. \\
    \code{pull\_requests} & Number, title, author, head and base commits, and open, closed, or merged. Unique per repository and number. \\
    \code{reviews} & One pipeline run. Pull request and repository are null for a local MCP review; \code{user\_id} is who asked. Stores provider, model, token counts, duration, summary, and error. \\
    \code{findings} & One comment: file, lines, severity, category, source (\code{llm} or \code{semgrep}), message, suggestion, confidence, and whether it was posted. \\
    \code{review\_events} & One row per pipeline stage, which is the timeline on the review page. \\
    \code{code\_chunks} & A symbol or a line window, its text, a content hash, and a vector. An HNSW index serves cosine search. \\
    \code{api\_keys} & A prefix and a hash. The raw key is returned once at creation and is not stored. \\
    \code{llm\_settings} & One row per user: provider keys and model choices. A null column means ``use the server default'', subject to the admin rule in Chapter~\ref{chap:architecture}. \\
    \code{github\_app} & The single App created from the dashboard (id, slug, private key, webhook secret, and the admin user). Environment variables are only a fallback. \\
    \code{notifications} & ``Review finished'', ``indexing failed'', and the other worker messages, with an optional dashboard link. \\
    \code{webhook\_deliveries} & Primary key is the GitHub delivery id. \\
    \bottomrule
  \end{tabular}
\end{table}

\section{Embeddings}

The vector column width is the \code{EMBEDDING\_DIM} used at migration time, 768 by default. It has to match the model. \code{repositories.embedding\_model} stores the \code{provider:model} pair that built the index. Indexing with a different embedder deletes the old chunks and rebuilds. Search and review context refuse an index that was built with another embedder, and the dashboard asks for a re-index.

Chunks are unique by content hash, so a later index does not re-embed a function that did not change.

\section{What is deliberately not logged}

The worker writes short stage messages to \code{review\_events}. Those messages do not include source code. The chunk table does store chunk text, because search and review context need it. The clone used during indexing is deleted when the job ends. API keys are stored as hashes. GitHub installation tokens are short-lived and are not written to a table.
